\section{Теоретические сведения}
\label{sec:theory}

\subsection{Линейное кодирование и восстановление}
\label{sec:theory-code}

Линейная рециркуляционная сеть является двухслойной автоассоциативной сетью. Для вектора-строки $x_k$:
\begin{equation}
	y_k = x_k W, \qquad \hat{x}_k = y_k W',
	\label{eq:forward}
\end{equation}
где $W \in \mathbb{R}^{n \times p}$, $W' \in \mathbb{R}^{p \times n}$, $0<p<n$. При $p<n$ скрытое представление $y_k$ содержит меньше компонент, чем вход, поэтому первый слой выполняет сжатие, а второй~--- восстановление. Оба слоя линейны, смещения и нелинейные активации отсутствуют~\cite{haykin}.

Для всей выборки сразу:
\begin{equation}
	Y = XW, \qquad \hat{X} = YW',
\end{equation}
где $X \in \mathbb{R}^{Q \times n}$, $Y \in \mathbb{R}^{Q \times p}$, $\hat{X} \in \mathbb{R}^{Q \times n}$.

При перекрытии блоков один пиксель входит в несколько блоков. При сборке изображения значения перекрывающихся блоков усредняются, что уменьшает видимость швов на границах.

\subsection{Нормализация и денормализация пикселей}
\label{sec:theory-norm}

Исходные значения каналов лежат в диапазоне $[0,C_{\max}]$, $C_{\max}=255$. Нормировка~\eqref{eq:norm} переводит их в $[-1,\,1]$. После работы сети выполняется обратное преобразование с ограничением диапазона:
\begin{equation}
	C' = \frac{C_{\max}\bigl(\min(1,\max(-1,\hat{c}))+1\bigr)}{2}.
	\label{eq:denorm}
\end{equation}
Ограничение необходимо, поскольку линейный выход сети может выйти за пределы $[-1,\,1]$. Ошибка нормированных блоков и ошибка готового BMP после округления и сборки вычисляются отдельно.

\subsection{Вывод градиентов}
\label{sec:theory-grad}

Запишем прямой проход покомпонентно:
\begin{equation}
	y_{k,j} = \sum_{i=1}^{n} x_{k,i} W_{ij}, \qquad
	\hat{x}_{k,i} = \sum_{j=1}^{p} y_{k,j} W'_{ji}.
	\label{eq:forward-comp}
\end{equation}
Ошибка одного образа:
\begin{equation}
	E_k = \frac{1}{2}\sum_{i=1}^{n}\bigl(\hat{x}_{k,i}-x_{k,i}\bigr)^2
	    = \frac{1}{2}\sum_{i=1}^{n} e_{k,i}^2 .
	\label{eq:ek}
\end{equation}

По правилу цепочки:
\begin{equation}
	\frac{\partial E_k}{\partial W'_{ji}}
	  = \frac{\partial E_k}{\partial \hat{x}_{k,i}}
	    \cdot \frac{\partial \hat{x}_{k,i}}{\partial W'_{ji}}
	  = e_{k,i}\,y_{k,j}.
	\label{eq:grad-wb}
\end{equation}
Следовательно, градиент по декодеру для одного образа есть внешнее произведение
\begin{equation}
	\nabla_{W'}E_k = y_k^{\mathsf{T}} e_k .
\end{equation}

Для кодера сначала перенесём ошибку в скрытый слой:
\begin{equation}
	\delta_{k,j} = \frac{\partial E_k}{\partial y_{k,j}}
	  = \sum_{i=1}^{n} \frac{\partial E_k}{\partial \hat{x}_{k,i}}
	    \cdot \frac{\partial \hat{x}_{k,i}}{\partial y_{k,j}}
	  = \sum_{i=1}^{n} e_{k,i} W'_{ji}
	  = \bigl(e_k W'^{\mathsf{T}}\bigr)_j .
	\label{eq:delta}
\end{equation}
Тогда
\begin{equation}
	\frac{\partial E_k}{\partial W_{ij}}
	  = \frac{\partial E_k}{\partial y_{k,j}}
	    \cdot \frac{\partial y_{k,j}}{\partial W_{ij}}
	  = \delta_{k,j}\,x_{k,i},
	\qquad
	\nabla_W E_k = x_k^{\mathsf{T}} \delta_k .
	\label{eq:grad-wf}
\end{equation}

В матричной форме для одного образа:
\begin{equation}
	\nabla_{W'}E_k = y_k^{\mathsf{T}} e_k, \qquad
	\nabla_W E_k = x_k^{\mathsf{T}} \bigl(e_k W'^{\mathsf{T}}\bigr).
	\label{eq:grads}
\end{equation}

\subsection{Дельта-правило и адаптивный коэффициент обучения}
\label{sec:theory-delta}

Обучение выполняется стохастическим градиентным спуском: веса обновляются после каждого образа в порядке, который на каждой эпохе перемешивается. Для варианта 2 используется адаптивный шаг~\eqref{eq:alpha}. С учётом~\eqref{eq:grads}:
\begin{equation}
	W' \leftarrow W' - \alpha_k\, y_k^{\mathsf{T}} e_k,
	\qquad
	W \leftarrow W - \alpha_k\, x_k^{\mathsf{T}} \delta_k .
	\label{eq:update}
\end{equation}

Нормировка весов после обновления \textbf{не выполняется}: матрицы $W$ и $W'$ сохраняются в том масштабе, который получился в ходе обучения. Это принципиальное отличие варианта 2 от варианта 11.

Адаптивный шаг~\eqref{eq:alpha} нормирует поправку полной энергией входа и скрытого представления. Если все значения образа умножить на постоянный множитель, числитель и знаменатель изменятся согласованно, поэтому направление и относительная величина обновления сохраняются. Параметр $\beta$ задаёт компромисс между устойчивостью и скоростью: при малых $\beta$ обучение медленное, при слишком больших~--- возможно расхождение.

Для локального квадратичного приближения ошибки в окрестности минимума
\begin{equation}
	E(w) \approx \frac{1}{2}(w-w^*)^{\mathsf{T}} H (w-w^*),
\end{equation}
где $H$~--- симметричная матрица Гессе. Тогда отклонение от минимума эволюционирует как
\begin{equation}
	w(t+1)-w^* = \bigl(I-\alpha H\bigr)\bigl(w(t)-w^*\bigr).
\end{equation}
По спектральной теореме для симметричных матриц~\cite{gantmacher} в собственном базисе $H=V\Lambda V^{\mathsf{T}}$ компоненты затухают независимо:
\begin{equation}
	z_i(t+1) = (1-\alpha\lambda_i)z_i(t), \qquad i=1,\ldots,n.
	\label{eq:eigen}
\end{equation}
Условие устойчивости $|1-\alpha\lambda_i|<1$ для всех $i$ даёт
\begin{equation}
	0 < \alpha < \frac{2}{\lambda_{\max}} .
	\label{eq:stability}
\end{equation}
Для линейного слоя с квадратичной ошибкой матрица кривизны пропорциональна $X^{\mathsf{T}}X$, поэтому $\lambda_{\max}>0$ при ненулевых вариациях данных. Адаптивный шаг~\eqref{eq:alpha} является практической оценкой допустимого масштаба; условие~\eqref{eq:stability} объясняет, почему при увеличении $\beta$ сверх критического значения метод расходится.

\subsection{Ошибка и критерий остановки}
\label{sec:theory-stop}

После каждой эпохи вычисляется суммарная ошибка~\eqref{eq:loss} по всем образам с текущими весами. Обучение прекращается при $E \le \varepsilon$ либо при исчерпании лимита $T_{\max}$. Ошибки прошлых эпох не накапливаются.

Средняя ошибка нормированных блоков связана с $E$:
\begin{equation}
	\mathrm{MSE}_n = \frac{2E}{Q\,n}.
	\label{eq:mse-n}
\end{equation}
Отдельно вычисляется $\mathrm{MSE}$ готового BMP в шкале $0\ldots255$:
\begin{equation}
	\mathrm{MSE}_{\mathrm{pix}} = \frac{1}{h\,w\,|S|}
	  \sum_{i,j,c}\bigl(C_{ijc}-C'_{ijc}\bigr)^2,
	\qquad
	\mathrm{PSNR} = 10\lg\frac{255^2}{\mathrm{MSE}_{\mathrm{pix}}}.
	\label{eq:psnr}
\end{equation}
Эти величины не смешиваются: первая характеризует ошибку до ограничения и сборки, вторая~--- качество итогового изображения.

\subsection{Коэффициент сжатия}
\label{sec:theory-comp}

По методическим указаниям коэффициент сжатия равен
\begin{equation}
	Z(Q,n,p) = \frac{Q\,n}{2+(n+Q)\,p}.
	\label{eq:Z-theory}
\end{equation}
Числитель $Qn$~--- объём исходных данных выборки блоков. В знаменателе $Qp$~--- объём скрытых кодов, $np$~--- объём общей матрицы декодера, а слагаемое $2$ соответствует двум служебным размерностям $n$ и $p$. Формула предполагает одинаковую разрядность исходных и сохранённых значений.

В программе исходные каналы хранятся в \texttt{uint8} (1 байт), а коды и декодер~--- в \texttt{float64} (8 байт). Поэтому дополнительно рассчитывается байтовый коэффициент
\begin{equation}
	K = \frac{Q\,n \cdot 1}{Q\,p \cdot 8 + p\,n \cdot 8 + 8}.
	\label{eq:K-byte}
\end{equation}
Здесь $8$ в конце~--- два служебных значения \texttt{uint32}. Величина $K$ может быть меньше единицы даже при $Z>1$: это означает, что переход от \texttt{uint8} к \texttt{float64} увеличивает объём, и коэффициент сжатия сам по себе ещё не гарантирует экономии байтов. В отчёте приводятся обе величины: $Z$~--- по методике, $K$~--- как уточнённая байтовая оценка.
